\subsection{批内参与方区分目标}

第\ref{sec:identity-problem}节将K-Waay的不同参与方条件与消息、发送实例、槽位和批次坐标区分开来。现将该批内关系形式化为分析目标。设$B=(E_1,\ldots,E_n)$为有序批次，$\mathsf{party}(E_i)$表示与条目$E_i$关联的建模协议主体，定义

\begin{equation}
  \distinctparty(B)
  \mathrel{\triangleq}
  \forall\,1 \leq i < j \leq n,\;
  \mathsf{party}(E_i)\neq\mathsf{party}(E_j).
  \label{eq:distinct-party-objective}
\end{equation}

因此，同一已接纳批次中任意两个不同位置都必须投影到不同的建模参与方。该目标是参与方坐标上的关系，不能由比较消息内容、发送实例标识或槽位索引得到。

式\eqref{eq:distinct-party-objective}的作用域限于一个批次。它既不要求某一参与方在整个协议、不同批次或不同会话中只出现一次，也不要求消息或发送实例全局唯一。若某个组合不满足该式，下文仅称其违反参与方级批次不变量；这种违反本身不表示密码材料畸形、认证失败或底层密码原语被攻破。

这一不变量是对批次组成的结构谓词，而非对单个条目的密码真实性谓词。它规定哪些参与方投影可以共存于一个批次，但不把参与方投影的获取方式和检查机制写入目标本身。定义针对任意有限$n$给出，以使目标关系独立于模型界限；机器分析固定$n=2$，此时目标化为$\mathsf{party}(E_1)\neq\mathsf{party}(E_2)$。当前证据不构成任意批大小的完整协议证明，也不预设违反该不变量必然产生接收侧可观察后果。

\subsection{对手能力与批次组合边界}

威胁模型采用符号批次组合对手（symbolic batch-composition adversary）。标准Dolev--Yao网络能力使对手能够获知、保留、重放并重新组合网络上暴露的符号项；在此基础上，模型另行假定对手能够在候选发送方条目已经可用之后，选择条目、安排其槽位，并把所得组合提交给接纳步骤。该能力位于候选条目可用性与批次接纳决定之间，只涉及输入组合，不授予对手创建诚实参与方内部状态的能力。

这一组合能力是一项分析假设，不能直接解释为真实网络对手已经控制部署中的批次构造器。具体系统可能由客户端、服务器、调用者或其他组件构造批次；若要把模型能力提升为部署结论，还需说明相应接口、缓冲与调度机制。本文没有提供这种实现层证据。

威胁模型区分两类重复输入。第一类是把同一暴露元组提交两次，即$E_1=E_2$。第二类是选取两个不同的发送来源元组

\begin{equation}
  E_1=(A,\mathit{oid}_1,m_1),
  \qquad
  E_2=(A,\mathit{oid}_2,m_2),
  \qquad
  E_1\neq E_2,
  \label{eq:repeated-party-capability}
\end{equation}

但二者具有相同的参与方投影$A$。在原型的新鲜值假设下，第二类元组的发送实例与消息坐标均不同。这里所需的是批次组合能力，而非身份伪造：对手无需冒充或虚构另一个建模参与方。

对于涉及已接收分量的结论，发送方与接收方之间的匹配作用于完整建模来源元组$(A,\mathit{oid},m)$，而非仅作用于$A$。对手可以在语法上重组从网络获知的值，但只有该完整元组确由建模发送生命周期产生时，重组结果才具有匹配来源；参与方坐标相等本身不提供来源证据。所分析的组合也不依赖密钥恢复、KEM或签名机制被攻破，亦不依赖伪造模型省略的密码材料。这些排除项只界定抽象范围，并不构成对被省略机制的安全证明。

\subsection{建模假设}

每次建模的发送发生均新鲜生成发送实例坐标$\mathit{oid}$和消息$m$。因此，两个不同的诚实发送来源在模型中具有不同的实例和消息坐标；同一参与方仍可触发多个发送发生。新鲜性用于控制三个模型的共同生命周期，不应被解释为真实协议消息不存在碰撞、重传或复用行为的独立结论。

来源匹配采用理想化精确关系：\receiveraccept 与先前的发送来源按完整$(A,\mathit{oid},m)$元组对应。仅有$A$相同不足以建立来源，来源元组中的值被重新拼接也不足以得到匹配。该假设使分析能够分离批次组合与来源伪造，但没有建模K-Waay中实现这一关联所需的全部密码计算。

机器模型固定为两个槽位，并以顺序生命周期处理已接纳分量。相关事件还被限定在同一批次标识和接收方上下文$(\mathit{bid},\mathit{rst})$内；性质不跨批次推断全局唯一性。两槽约束是验证边界，而非式\eqref{eq:distinct-party-objective}的定义边界。模型比较由此针对相同的发送、收集、来源匹配和处理骨架，只改变接纳谓词所约束的身份坐标。

\subsection{证据边界与非目标}

直接形式化对象包括产生来源元组所需的发送生命周期、符号网络上的批次组合、批次接纳或拒绝、发送来源关系以及分量级\receiveraccept 事件。\receiveraccept 是批处理接纳抽象中的模型局部事件，并非K-Waay原安全模型定义的协议级接受事件；引入该事件本身不预设任何验证结果。

机器检查证据止于\receiveraccept。当前模型没有展开完整的KEM、数字签名和KDF计算，也没有妥协行为、棘轮状态、下游密钥接口或应用消费者；模型中不存在正式的\textsc{Key}或\textsc{Test}对象，不输出供上层使用的会话密钥，也不包含应用状态安装事件。因此，结果不直接建立机密性、密钥不可区分性、完整协议认证失败或重复安装后果。

将接收事件连接到上层状态需要额外的组合假设，本文不作此断言。当前工作同样没有给出从完整K-Waay执行到抽象模型的精化定理，没有验证任意有限批大小，也不评价某个部署实现是否可被利用。其证据范围是：在上述固定两槽、精确来源和批次上下文假设下，比较不同身份坐标上的接纳限制所保持的关系。
